% Generated draft: requires mathematical and source review.
\documentclass[11pt]{article}
\usepackage[margin=25mm]{geometry}
\usepackage{fontspec}
\setmainfont{Latin Modern Roman}
\usepackage{amsmath,amssymb,mathtools,unicode-math}
\setmathfont{Latin Modern Math}
\usepackage{enumitem,needspace,xurl,hyperref}
\setlength{\parindent}{0pt}
\setlength{\parskip}{5pt}
\setlength{\emergencystretch}{3em}
\newcommand{\N}{\mathbb{N}}
\newcommand{\Z}{\mathbb{Z}}
\newcommand{\Q}{\mathbb{Q}}
\newcommand{\R}{\mathbb{R}}
\newcommand{\C}{\mathbb{C}}
\newcommand{\F}{\mathbb{F}}
\newcommand{\A}{\mathbb{A}}
\newcommand{\PP}{\mathbb{P}}
\newcommand{\E}{\mathbb{E}}
\newcommand{\eps}{\varepsilon}
\newcommand{\dd}{\,\mathrm{d}}
\begin{document}
\section*{UnsolvedMath ID: 2301005}
\textbf{Source problem code:} catalogue entry
\subsection{Definitions and mathematical statement}
For a nonconstant entire function $f$, put
\[
T(r,f)=\frac1{2\pi}\int_0^{2\pi}\log^+|f(re^{i\theta})|\,d\theta,
\qquad
\rho=\limsup_{r\to\infty}\frac{\log T(r,f)}{\log r},\quad
\lambda=\liminf_{r\to\infty}\frac{\log T(r,f)}{\log r}.
\]
Here $\log^+t=\max(\log t,0)$, $\rho$ is the order, and $\lambda$ is the lower order. Let $N(r,a;f)$ be the integrated counting function of the $a$-points, with multiplicity, and define the Nevanlinna deficiency by
\[
\delta(a,f)=1-\limsup_{r\to\infty}\frac{N(r,a;f)}{T(r,f)}.
\]
For $a=\infty$, count poles; hence $\delta(\infty,f)=1$ for entire $f$. To interpret a total near $2$, use
\[
D(f)=\sum_{a\in\widehat{\mathbb C}}\delta(a,f),
\]
including infinity. The record does not explicitly state this summation convention; summing only over finite values is a different assertion.

\textbf{Problem.} For entire $f$ of finite order $\rho$, characterize how $D(f)$ can be nearly $2$. In particular, in the source's regime
\[
D(f)>2-\varepsilon(\lambda),
\]
can there also be finitely many or infinitely many small positive deficiencies, in addition to the large deficiencies described below?

The source attributes to Pfluger the statement that $D(f)=2$ forces $\rho=q$ for a positive integer $q$, $\lambda=\rho$, and every positive deficiency to be an integral multiple of $1/q$. It attributes to Edrei and Fuchs an approximate version under the displayed near-maximal inequality: certain large deficiencies are near positive integral multiples of $1/q$ and their sum is near $2$. The positive function $\varepsilon(\lambda)$, the selection of $q$, and the thresholds and error bounds for nearly, large, and small are not provided. Thus the request about additional small deficient values is retained without inventing a quantitative cutoff. The supplied summary describes the problem as open in the 2018 edition, with no later resolution located in its search.
\subsection{Short English statement}
When an entire function nearly exhausts the total deficiency bound, can it have additional small deficiencies, finitely or infinitely many, alongside the large nearly quantized ones? The source’s quantitative thresholds and theorem references are incomplete.
\subsection{Sources}
\begin{itemize}
\item Hayman, \emph{Research Problems in Function Theory} (2018), Problem 1.5, \url{https://arxiv.org/abs/1809.07200}; source-TeX extraction, accessed 2026-07-29. The record retains the 2018 update for status triage.
\item Pfluger, equality-case theorem, with an incomplete see-Hayman reference in the source; Edrei and Fuchs, near-equality results, with two missing citation entries in the extraction. No titles, dates, or theorem locations for these results are supplied.
\item Local registry \texttt{data/unsolvedmath/problems.json}, permanent ID $2301005$, source code catalogue entry, \emph{Research Problems in Function Theory --- Problem 1.5}. The record labels this open and OPEN-TRIAGE, with NEEDS\_REVIEW source-status and rights qualifications.
\end{itemize}
\end{document}
